\documentclass[../../main.tex]{subfiles}

\begin{document}

\section{基本概念}\label{section:ch01-基本概念}

我们将在本节介绍集合论的思想，并建立基本术语和记号。我们还将讨论初等逻辑学的某些内容，根据我们的经验，这些内容容易引起混淆。

\subsection*{基本记号}

通常我们用大写字母 $A$，$B$，$\cdots$ 表示集合（set），用小写字母 $a$，$b$，$\cdots$ 表示属于集合的成员（object）或元素（element）。集合有时简称为集（set），元素有时简称为元（element）或点（point）。如果成员 $a$ 属于集合 $A$，就记作
\begin{align*}
a \in A.
\end{align*}
如果 $a$ 不属于 $A$，就记作
\begin{align*}
a \notin A.
\end{align*}

本书中用到的等号 $=$ 是指逻辑上的同一。当我们写 $a = b$ 时，就意味着“$a$”和“$b$”是同一个成员的两个符号。这就像在算术中写 $\dfrac{2}{4} = \dfrac{1}{2}$ 一样。类似地，$A = B$ 就是说“$A$”和“$B$”是同一个集合的两个符号，也就是说 $A$ 与 $B$ 含有完全相同的成员。如果 $a$ 与 $b$ 是不同的成员，就写作 $a \neq b$；如果 $A$ 与 $B$ 是不同的集合，就写作 $A \neq B$。例如，设 $A$ 为所有非负实数的集合，$B$ 为所有正实数的集合，由于数 $0$ 属于 $A$ 而不属于 $B$，所以 $A \neq B$。

如果 $A$ 的每一个元素都是 $B$ 的元素，就说 $A$ 是 $B$ 的子集（subset），记作
\begin{align*}
A \subset B.
\end{align*}
这个定义中并没有要求 $A$ 不能等于 $B$。事实上，如果 $A = B$，那么 $A \subset B$ 与 $B \supset A$ 都成立。当 $A \subset B$ 并且 $A$ 不等于 $B$ 时，称 $A$ 为 $B$ 的真子集（proper subset），记作
\begin{align*}
A \subsetneqq B.
\end{align*}
两集合间的 $\subset$ 和 $\subsetneqq$ 关系分别称之为包含（inclusion）关系和真包含（proper inclusion）关系。若 $A \subset B$，则也可写作 $B \supset A$，读作“$B$ 包含 $A$”。

怎样来描述一个集合呢？如果它只含有为数不多的元素，那么可以把集合中的成员都列出来，写作“$A$ 是由元素 $a$、$b$ 和 $c$ 组成的集合”。使用符号来表达就是
\begin{align*}
A = \{a, b, c\},
\end{align*}
这里的花括号用来把列举出来的所有元素包在一起。

刻画一个由某些成员构成的集合，最常用的方法是取一个集合 $A$ 和 $A$ 中的成员可能具有也可能不具有的某种性质，然后用 $A$ 中具有这种性质的所有元素来组成集合。例如，可以取实数集，并且由所有偶数组成其子集 $B$。使用符号将这句话写成
\begin{align*}
B = \{x \mid x \text{ 是偶数}\}.
\end{align*}
在这里，花括号表示“$\cdots\cdots$ 组成的集合”这个词组，竖线表示“使得”这个词，整个式子读作“$B$ 是所有使得 $x$ 为偶数的那些 $x$ 组成的集合”。

\subsection*{集合的“并”与“或”的含义}

给定两个集合 $A$ 和 $B$，由 $A$ 中所有元素及 $B$ 中所有元素可以组成一个集合，这个集合称为 $A$ 与 $B$ 的并（union）或并集，记作 $A \cup B$。正式的定义是
\begin{align}
A \cup B = \{x \mid x \in A \text{ 或 } x \in B\}.\label{align:ch01-并定义}
\end{align}
现在我们必须停下来，看一看“$x \in A$ 或 $x \in B$”这句话究竟意味着什么。

在日常用语中，“或”这个词是含糊的。有时“$P$ 或 $Q$”这句话意味着“$P$ 或 $Q$，或者既 $P$ 又 $Q$”，有时又意味着“$P$ 或 $Q$，但不是既 $P$ 又 $Q$”。通常这要从文章的上下文才能知道究竟指的是哪一种。例如，我对两个学生说：

\begin{quote}
“Smith 小姐，每一个选修这门课的学生，或者学过线性代数，或者学过数学分析。”

“Jones 先生，你这门课程的期末考试，或者不低于 $70$ 分，或者不及格。”
\end{quote}

从上下文看，Smith 小姐完全知道我说的意思是“每一个人要么学过线性代数，要么学过数学分析，要么两门课都学过”，Jones 先生也明白我说的是“或者他至少得 $70$ 分，或者他不及格，两者仅取其一”。因为，如果两者同时成为事实的话，Jones 先生就太不幸了。

数学中不能容许这种含糊。自始至终只能承认它的一种含义，否则就要引起混乱。因此，数学家们同意在第一种意义下使用“或”这个词，这样，“$P$ 或 $Q$”这句话总是指“$P$ 或 $Q$，或者既 $P$ 又 $Q$”。如果要指“$P$ 或 $Q$，但不是既 $P$ 又 $Q$”，就必须明确地加上短语“但不是既 $P$ 又 $Q$”。按照这种解释，定义 $A \cup B$ 的式子就清楚了，它表明 $A \cup B$ 是由所有属于 $A$，或者属于 $B$，或者既属于 $A$ 又属于 $B$ 的元素 $x$ 组成的集合。

\subsection*{集合的交、空集以及“若$\cdots$，则$\cdots$”的含义}

给定两个集合 $A$ 和 $B$，还可以用另一种方法组成一个集合，就是取 $A$ 与 $B$ 的公共部分。这个集合称为 $A$ 与 $B$ 的交（intersection）或交集，记作 $A \cap B$。正式的定义是
\begin{align}
A \cap B = \{x \mid x \in A \text{ 和 } x \in B\}.\label{align:ch01-交定义}
\end{align}
与 $A \cup B$ 的定义一样，这里也有一个麻烦。当然，它与前面提及的问题有所区别，是另一类麻烦。现在的问题不在于“和”这个词的含义，而在于当 $A$ 与 $B$ 没有公共元素时，记号 $A \cap B$ 意味着什么？

为了应付这种偶然的情形，需要作一个特殊的约定。我们引进一个称为空集（empty set）的特殊的集合，记成 $\varnothing$，设想成“没有元素的集合”。

使用这种约定，$A$ 与 $B$ 没有公共元素这句话就记作
\begin{align*}
A \cap B = \varnothing.
\end{align*}
这时，也说 $A$ 与 $B$ 无交（disjoint）。

一些学生对于“空集”的概念感到困惑不解。他们说：“你怎么能够找到一个集合，而它什么也没有呢？”这和多年以前人们首次引进数 $0$ 时所涉及的问题是一样的。

空集仅仅是一个约定，数学中完全可以不要它。但是，有这种约定就比较方便，它能使定理的叙述和证明变得更加简洁。例如，不用这个约定，那就必须在使用记号 $A \cap B$ 之前，证明两个集合 $A$ 和 $B$ 有公共元素。类似地，对于记号
\begin{align*}
C = \{x \mid x \in A \text{ 并且 } x \text{ 有某性质}\},
\end{align*}
当 $A$ 中没有具有给定性质的元素 $x$ 时，它就无法使用。这种时候如果认为 $A \cap B$ 和 $C$ 是空集就方便多了。

由于空集 $\varnothing$ 只不过是一个约定，所以要对前面引进的概念作出与空集相关的约定。空集既然被看成“没有元素的集合”，显然可以约定：对于每一个成员 $x$，关系式 $x \in \varnothing$ 不成立。类似地，由“并”和“交”的定义可见，对于任意集合 $A$，有
\begin{align*}
A \cup \varnothing = A \quad \text{和} \quad A \cap \varnothing = \varnothing.
\end{align*}

包含关系也有点小麻烦。给定一个集合 $A$，可以认为 $\varnothing \subset A$ 吗？我们必须再次注意数学家们运用语言的方式。表达式 $\varnothing \subset A$ 其实是语句“每一属于空集的元素都属于集合 $A$”的缩写。或者更正式地说：“对于每一个成员 $x$，若 $x$ 属于空集，则 $x$ 必属于集合 $A$。”

这种说法对不对呢？有人说“对”，也有人说“不对”。这个问题很难加以论证，只能予以约定。它是“若 $P$，则 $Q$”式的论断。在日常用语中，“若$\cdots$，则$\cdots$”结构的含义是什么？通常指的是：若 $P$ 为真，则 $Q$ 也为真。有时候这就是它所包含的全部内容。有时候它还有另外的含义，即：若 $P$ 不真，则 $Q$ 必不真。通常我们可以根据上下文来判定其真正的含义。

这件事与使用“或”这个词时所造成的混乱很类似。我们还用前面那个关于 Smith 小姐和 Jones 先生的例子来说明其中的歧义。假定我说：

\begin{quote}
“Smith 小姐，每一个选修这门课的学生，若没有学过线性代数，则必须学过数学分析。”

“Jones 先生，若你期末考试低于 $70$ 分，你这门课就不及格。”
\end{quote}

从上下文看，Smith 小姐知道，如果一个学生要学习这门课程，他没有学过线性代数，则必须学过数学分析；但是，如果学过线性代数，那么，他可以学过数学分析，也可以没有学过数学分析。Jones 先生也明白，如果他的考试成绩低于 $70$ 分，他这门课就不及格；如果他至少得了 $70$ 分，那就及格了。

再次申明：数学不允许这种模棱两可，对于它的含义只能作出一种选择。数学家们往往同意对于“若$\cdots$，则$\cdots$”句式采取第一种解释，所以“若 $P$，则 $Q$。”式论断的含义是：若 $P$ 为真，则 $Q$ 也真；若 $P$ 不真，则 $Q$ 可以真，也可以不真。

举一个例子，考虑下面关于实数的一个论断：
\begin{align*}
\text{若 } x > 0, \quad \text{则 } x^{3} \neq 0.
\end{align*}
这是“若 $P$，则 $Q$。”式的一个论断，其中 $P$ 代表短语“$x > 0$”（称为论断假设（hypothesis）），$Q$ 代表短语“$x^{3} \neq 0$”（称为论断结论（conclusion））。它是真论断，因为在任何情况下，只要假设 $x > 0$ 成立，则结论 $x^{3} \neq 0$ 必成立。

再看一个关于实数的真论断：
\begin{align*}
\text{若 } x^{2} < 0, \quad \text{则 } x = 23;
\end{align*}
不论何时，只要假设成立，则结论总是成立的。当然，这里的假设在任何情况下也不会成立。这一类论断有时称为虚真论断（vacuously true）。

现在再回到空集和包含关系上，我们看到对于任何集合 $A$，包含关系 $\varnothing \subset A$ 必成立。$\varnothing \subset A$ 就等于说：“若 $x \in \varnothing$，则 $x \in A$。”而这是一个虚真论断。

\subsection*{逆否论断与逆论断}

对于“若$\cdots$，则$\cdots$”结构的讨论，使我们可以研究初等逻辑学中另外一些更为复杂的问题，这就是一个论断与它的逆否论断、逆论断之间的关系。

给定一个“若 $P$，则 $Q$。”式的论断，它的逆否论断（contrapositive）定义为“若 $Q$ 不真，则 $P$ 不真。”例如论断
\begin{align*}
\text{若 } x > 0, \quad \text{则 } x^{3} \neq 0,
\end{align*}
其逆否论断是
\begin{align*}
\text{若 } x^{3} = 0, \quad \text{则 } x > 0 \text{ 不成立}.
\end{align*}
我们发现，上述论断和它的逆否论断同时为真。类似地，论断
\begin{align*}
\text{若 } x^{2} < 0, \quad \text{则 } x = 23,
\end{align*}
其逆否论断是
\begin{align*}
\text{若 } x \neq 23, \quad \text{则 } x^{2} < 0 \text{ 不成立}.
\end{align*}
就实数而言，这两个论断也都正确。

这些例子可能会使你感到，在一个论断与它的逆否论断之间有着某种关系。事实正是如此，它们是同一件事情的两种不同的说法。任何一个为真的充分必要条件是另一个为真，它们是逻辑等价的。

这个事实不难证明。我们首先引进一个记号。将“若 $P$，则 $Q$。”式的论断简单地记为
\begin{align*}
P \Longrightarrow Q,
\end{align*}
读作“$P$ 蕴涵 $Q$”，则其逆否论断可以写成
\begin{align*}
(\text{非 } Q) \Longrightarrow (\text{非 } P),
\end{align*}
其中“非 $Q$”表示短语“$Q$ 不真”。

说明论断“$P \Longrightarrow Q$”不成立的唯一办法是：当假设 $P$ 为真时，结论 $Q$ 不真。如若不然，论断就成立了。类似地，说明论断 $(\text{非 } Q) \Longrightarrow (\text{非 } P)$ 不成立的唯一办法是：当假设“非 $Q$”为真时，结论“非 $P$”不真。这等于说 $Q$ 不真而 $P$ 真，这又恰好描述了 $P \Longrightarrow Q$ 不成立。因此，这两个论断或同时为真，或同时为不真，它们逻辑等价。所以我们认为，证明论断“非 $Q \Longrightarrow$ 非 $P$”，就是证明了论断“$P \Longrightarrow Q$”。

由论断 $P \Longrightarrow Q$ 还可以得到另一个论断，那就是
\begin{align*}
Q \Longrightarrow P,
\end{align*}
称为 $P \Longrightarrow Q$ 的逆（converse）。必须注意一个论断的逆和它的逆否之间的不同。论断与其逆否论断逻辑等价，而论断为真，却完全不能保证其逆论断是否正确。例如，真论断
\begin{align*}
\text{若 } x > 0, \text{则 } x^{3} \neq 0,
\end{align*}
其逆论断
\begin{align*}
\text{若 } x^{3} \neq 0, \text{则 } x > 0,
\end{align*}
就是不真的。同样，真论断
\begin{align*}
\text{若 } x^{2} < 0, \text{则 } x = 23,
\end{align*}
其逆论断
\begin{align*}
\text{若 } x = 23, \text{则 } x^{2} < 0,
\end{align*}
也是非真论断。

如果恰好论断 $P \Longrightarrow Q$ 与其逆论断 $Q \Longrightarrow P$ 都为真，则将其记为
\begin{align*}
P \Longleftrightarrow Q,
\end{align*}
读作“$P$ 为真当且仅当 $Q$ 为真。”

\subsection*{否论断}

要想做出论断 $P \Longrightarrow Q$ 的逆否论断，必须知道如何做出论断“非 $P$”，称为 $P$ 的否定（negation）。这在大多数情形下并不困难，麻烦的是论断中包含有短语“对于任意”及“对于至少一个”。这类短语称为逻辑量词。

为了说明这一点，设 $X$ 是一个集合，$A$ 为 $X$ 的一个子集，$P$ 是关于 $X$ 中一般元素的一个论断。考虑下述论断
\begin{align}
\text{对于任意 } x \in A, \text{论断 } P \text{ 成立}.\label{align:ch01-全称论断}
\end{align}
如何做出它的否论断呢？我们把问题换成集合的语言。设 $B$ 表示 $X$ 中所有使论断 $P$ 成立的元素 $x$ 的集合，则论断\eqref{align:ch01-全称论断}正好是论断：$A$ 是 $B$ 的一个子集。其否论断是什么呢？显然应该是：$A$ 不是 $B$ 的子集，也就是说至少存在 $A$ 的一个元素，它不属于 $B$。再把它翻译成通常的语言就是
\begin{align*}
\text{对于至少一个 } x \in A, \text{论断 } P \text{ 不成立}.
\end{align*}
这就是论断\eqref{align:ch01-全称论断}的否论断，它用量词“对于至少一个”代替量词“对于任意”，用 $P$ 的否定代替 $P$。

正好把上述过程反过来做，就得到：对于论断
\begin{align*}
\text{对于至少一个 } x \in A, \text{论断 } Q \text{ 成立},
\end{align*}
其否论断是
\begin{align*}
\text{对于任意 } x \in A, \text{论断 } Q \text{ 不成立}.
\end{align*}

\subsection*{两个集合的差}

我们再来讨论有关集合的问题。在集合间还有另外一种常用运算，这就是两个集合的差（difference）或差集，记作 $A - B$，它由 $A$ 中所有不属于 $B$ 的元素组成。写成
\begin{align}
A - B = \{x \mid x \in A \text{ 和 } x \notin B\}.\label{align:ch01-差定义}
\end{align}
有时也称之为 $B$ 相对于 $A$ 的补（complement）或补集，或 $B$ 在 $A$ 中的补。集合的三种运算如\cref{figure:ch01-集合并交差}所示。

\par
\begin{figure}[htbp]
\centering
\includegraphics[width=0.85\textwidth]{ch01-fig1-1}
\caption{集合并、交、差的 Venn 图}\label{figure:ch01-集合并交差}
\end{figure}

\subsection*{集合论的法则}

对给定的一些集合施行集合运算，可以得到一些新的集合。就像代数中那样，我们也用括号表示运算的次序。例如 $A \cup (B \cap C)$ 表示集合 $A$ 与集合 $B \cap C$ 的并，$(A \cup B) \cap C$ 则表示集合 $A \cup B$ 与集合 $C$ 的交，两者完全不同，如\cref{figure:ch01-结合律}所示。

\par
\begin{figure}[htbp]
\centering
\includegraphics[width=0.7\textwidth]{ch01-fig1-2}
\caption{括号对运算次序的影响}\label{figure:ch01-结合律}
\end{figure}

有时按不同方式组合运算能得到相同的集合。这时我们便得到了集合论的一个法则。例如，对于任意集合 $A$，$B$，$C$，有以下公式
\begin{align}
A \cap (B \cup C) = (A \cap B) \cup (A \cap C).\label{align:ch01-第一分配律}
\end{align}
如\cref{figure:ch01-分配律}中所示，图中阴影部分表示上面的集合，你可以验证一下，它与你想像的是否一致。这个式子可以看成是关于运算 $\cap$ 和 $\cup$ 的一个“分配律”。

\par
\begin{figure}[htbp]
\centering
\includegraphics[width=0.35\textwidth]{ch01-fig1-3}
\caption{第一分配律的 Venn 图}\label{figure:ch01-分配律}
\end{figure}

集合论的另外两个法则是第二“分配律”
\begin{align}
A \cup (B \cap C) = (A \cup B) \cap (A \cup C),\label{align:ch01-第二分配律}
\end{align}
和 DeMorgan 定律
\begin{align}
A - (B \cup C) = (A - B) \cap (A - C),\label{align:ch01-DeMorgan并}
\end{align}
\begin{align}
A - (B \cap C) = (A - B) \cup (A - C).\label{align:ch01-DeMorgan交}
\end{align}
这些法则请读者自行验证。还可以举出集合论的另外一些法则，但是上面几个是最重要的。如果用语言叙述 DeMorgan 定律将更便于记忆，

\begin{quote}
并的补等于补的交。

交的补等于补的并。
\end{quote}

\subsection*{集合的族}

属于一个集合的成员可以是各式各样的。比如所有偶整数的集合，内布拉斯加州（Nebraska）所有蓝眼睛人的集合，世界上所有各种各样扑克牌的集合。这里面有些东西数学家并不感兴趣，可以不管它！但是，对于扑克牌的例子需要作一些说明，那就是集合的成员，其本身也可以是一个集合。一副牌是一个集合，它的成员是印了标准图案的单张扑克牌。世界上所有整副扑克牌的集合是以（单张扑克牌的）集合为成员的一个集合。

现在有了从已知集合得到新集合的另一个方法。给定集合 $A$，考虑其成员是 $A$ 的子集的集合。特别地，可以考虑 $A$ 的所有子集的集合，有时将其记作 $\mathcal{P}(A)$，称为 $A$ 的幂集（power set）（后面要说明这样称呼的理由）。

通常我们把一个以集合作为元素的集合称为族（collection），用 $\mathcal{A}$ 或 $\mathcal{B}$ 这样的花体字母来记。我们也常将“某些集合的族”说成集族（collection of sets）。这种记法能使我们在需要同时描述成员、成员的集合以及成员的族时，不致发生混淆。例如，我们用 $\mathcal{A}$ 表示世界上所有整副扑克牌的族，用大写字母 $A$ 记一副扑克牌，而用小写字母 $a$ 记一张扑克牌。

在这些记号同时出现时要格外小心。作为集合 $A$ 的一个元素的 $a$，与作为 $A$ 的子集的单点集 $\{a\}$ 是不同的。例如，若 $A$ 为集合 $\{a, b, c\}$，那么 $a \in A$，$\{a\} \subset A$，$\{a\} \in \mathcal{P}(A)$ 都是正确的，而 $\{a\} \in A$ 和 $a \subset A$ 则是错误的。

\subsection*{任意并与任意交}

我们已经定义了两个集合的并与交。当然没有理由只限于两个集合，同样可以讨论任意多个集合的并与交。

给定一个集族 $\mathcal{A}$，$\mathcal{A}$ 中元素的并（union）定义为
\begin{align}
\bigcup_{A \in \mathcal{A}} A = \{x \mid \text{对于至少一个 } A \in \mathcal{A}, x \in A\}.\label{align:ch01-任意并}
\end{align}
$\mathcal{A}$ 中元素的交（intersection）定义为
\begin{align}
\bigcap_{A \in \mathcal{A}} A = \{x \mid \text{对于任意 } A \in \mathcal{A}, x \in A\}.\label{align:ch01-任意交}
\end{align}

当 $\mathcal{A}$ 中有一个元素是空集时，这些定义没有问题。但是当 $\mathcal{A}$ 为空族时，要讲清楚这些定义的含义就有点麻烦了。按照定义的字面理解，任何元素 $x$ 都不满足 $\mathcal{A}$ 的成员的“并”所满足的性质。因此，当 $\mathcal{A}$ 为空族时，完全可以说
\begin{align*}
\bigcup_{A \in \mathcal{A}} A = \varnothing.
\end{align*}
另一方面，任意 $x$ 都（虚真地）满足 $\mathcal{A}$ 的成员的“交”所满足的性质。问题是，这个任意 $x$ 在哪个集合中呢？如果我们在一开始就给定一个包罗万象的大集合 $X$，并且始终只限于讨论 $X$ 的子集的话，那么，对于空族 $\mathcal{A}$，完全有理由令
\begin{align*}
\bigcap_{A \in \mathcal{A}} A = X.
\end{align*}
当然，不是所有的数学家都遵从这个约定。为了避免麻烦，我们将不定义空族 $\mathcal{A}$ 的交。

\subsection*{笛卡儿积}

再介绍一个由已知集合得到新集合的方法，它用到成员“有序偶对”的概念。在学习解析几何的时候，第一件事就是要确认，在平面上取定 $x$ 轴和 $y$ 轴后，平面上的任意一个点对应着唯一的一个有序实数对 $(x, y)$。（按照几何学中更为高深一点的说法，把平面定义为所有的有序实数对的集合将更为适宜些！）

将有序偶对的概念用在一般集合上。给定集合 $A$ 和 $B$，其笛卡儿积（Cartesian product）$A \times B$——笛卡儿积有时也称为积（product）——定义为所有有序偶对 $(a, b)$ 的集合，其中 $a$ 是 $A$ 的元素，$b$ 是 $B$ 的元素。记作
\begin{align}
A \times B = \{(a, b) \mid a \in A \text{ 和 } b \in B\}.\label{align:ch01-笛卡儿积}
\end{align}

这个定义是假定已经有了“有序偶对”的概念，它可以像“集合”那样，作为一个原始概念，也可以用前面讲过的集合运算来加以定义。其中一种表示为
\begin{align*}
(a, b) = \{\{a\}, \{a, b\}\};
\end{align*}
它将有序偶对 $(a, b)$ 定义为一个集族。如果 $a \neq b$，按照定义，$(a, b)$ 是两个集合的族，一个是单元素集，另一个是两元素集。有序偶对的第一个坐标是同时属于这两个集合的元素，第二个坐标是仅属于一个集合的元素。如果 $a = b$，这时 $\{a, b\} = \{a, a\} = \{a\}$，所以 $(a, b)$ 是仅一个集合 $\{a\}$ 的族。其第一、第二个坐标都是这个单点集的元素。

应当说大多数数学家更倾向于把有序偶对当作原始概念对待，而不将它作为集族。

我们对于记号作一点说明。因为很不凑巧，在数学中，记号 $(a, b)$ 已经有了两种完全不同的含义：一种含义是刚才讲的作为成员的一个有序偶对；另一种含义是我们在数学分析中早已熟悉的，即当 $a$，$b$ 为实数时，$(a, b)$ 表示包含所有满足不等式 $a < x < b$ 的所有数 $x$ 所构成的区间。由于它的含义可依行文而定，所以大多数情况不会有问题。在可能发生混淆时，我们就用另一个记号
\begin{align*}
a \times b
\end{align*}
来表示有序偶对 $(a, b)$。

\subsection*{习题}

\begin{exercise}\label{exercise:ch01-基本概念-1}
验证集合运算 $\cup$ 与 $\cap$ 的分配律及 DeMorgan 定律。
\end{exercise}

\begin{exercise}\label{exercise:ch01-基本概念-2}
对于所有集合 $A$，$B$，$C$，$D$，判定下面的哪些论断为真。如果相互蕴涵关系不成立，判定有哪一种蕴涵关系。如果等号不成立，判定用包含关系 $\subset$ 或 $\supset$ 代替等号时，相应的表述是否成立？
\begin{enumerate}[(a)]
\item $A \subset B$ 和 $A \subset C \Longleftrightarrow A \subset (B \cup C)$.
\item $A \subset B$ 或 $A \subset C \Longleftrightarrow A \subset (B \cup C)$.
\item $A \subset B$ 和 $A \subset C \Longleftrightarrow A \subset (B \cap C)$.
\item $A \subset B$ 或 $A \subset C \Longleftrightarrow A \subset (B \cap C)$.
\item $A - (A - B) = B$.
\item $A - (B - A) = A - B$.
\item $A \cap (B - C) = (A \cap B) - (A \cap C)$.
\item $A \cup (B - C) = (A \cup B) - (A \cup C)$.
\item $(A \cap B) \cup (A - B) = A$.
\item $A \subset C$ 并且 $B \subset D \Longrightarrow (A \times B) \subset (C \times D)$.
\item (j) 的逆。
\item (j) 的逆，假定 $A$ 和 $B$ 为非空集合。
\item $(A \times B) \cup (C \times D) = (A \cup C) \times (B \cup D)$.
\item $(A \times B) \cap (C \times D) = (A \cap C) \times (B \cap D)$.
\item $A \times (B - C) = (A \times B) - (A \times C)$.
\item $(A - B) \times (C - D) = (A \times C - B \times C) - (A \times D)$.
\item $(A \times B) - (C \times D) = (A - C) \times (B - D)$.
\end{enumerate}
\end{exercise}

\begin{exercise}\label{exercise:ch01-基本概念-3}
(a) 已知论断：“若 $x < 0$，则 $x^{2} - x > 0$。”试写出其逆否论断和逆论断。并判定三个论断中哪一个（如果有的话）是真论断。

(b) 将上述论断换成：“$x > 0$，则 $x^{2} - x > 0$。”考虑与 (a) 小题相仿的问题。
\end{exercise}

\begin{exercise}\label{exercise:ch01-基本概念-4}
设 $A$，$B$ 为实数集，写出下列每一个论断的否论断：
\begin{enumerate}[(a)]
\item 对于任何 $a \in A$，有 $a^{2} \in B$.
\item 存在一个 $a \in A$，有 $a^{2} \in B$.
\item 对于任何 $a \in A$，有 $a^{2} \notin B$.
\item 存在一个 $a \notin A$，有 $a^{2} \in B$.
\end{enumerate}
\end{exercise}

\begin{exercise}\label{exercise:ch01-基本概念-5}
设 $\mathcal{A}$ 是集的一个非空族。试判定下列每一个论断及其逆论断的真与假。
\begin{enumerate}[(a)]
\item $x \in \bigcup_{A \in \mathcal{A}} A \Longrightarrow$ 至少存在一个 $A \in \mathcal{A}$，使得 $x \in A$.
\item $x \in \bigcup_{A \in \mathcal{A}} A \Longrightarrow$ 对于任意 $A \in \mathcal{A}$，有 $x \in A$.
\item $x \in \bigcap_{A \in \mathcal{A}} A \Longrightarrow$ 至少存在一个 $A \in \mathcal{A}$，使得 $x \in A$.
\item $x \in \bigcap_{A \in \mathcal{A}} A \Longrightarrow$ 对于任意 $A \in \mathcal{A}$，有 $x \in A$.
\end{enumerate}
\end{exercise}

\begin{exercise}\label{exercise:ch01-基本概念-6}
写出习题\cref{exercise:ch01-基本概念-5}中每一个论断的逆否论断。
\end{exercise}

\begin{exercise}\label{exercise:ch01-基本概念-7}
已知集合 $A$，$B$，$C$，将下面每一个集合用集合 $A$，$B$，$C$ 及符号 $\cup$、$\cap$、$-$ 表示出来，
\begin{align*}
D = \{x \mid x \in A \text{ 并且 } (x \in B \text{ 或者 } x \in C)\},
\end{align*}
\begin{align*}
E = \{x \mid (x \in A \text{ 并且 } x \in B) \text{ 或者 } x \in C\},
\end{align*}
\begin{align*}
F = \{x \mid x \in A \text{ 并且 } (x \in B \Longrightarrow x \in C)\}.
\end{align*}
\end{exercise}

\begin{exercise}\label{exercise:ch01-基本概念-8}
若集合 $A$ 有两个元素，证明 $\mathcal{P}(A)$ 有四个元素。如果 $A$ 分别为单元素集、三元素集、空集，$\mathcal{P}(A)$ 又各有多少元素？为什么 $\mathcal{P}(A)$ 称为 $A$ 的幂集？
\end{exercise}

\begin{exercise}\label{exercise:ch01-基本概念-9}
对于任意“并”和任意“交”，陈述和证明 DeMorgan 定律。
\end{exercise}

\begin{exercise}\label{exercise:ch01-基本概念-10}
设 $\mathbb{R}$ 为实数集，试判定 $\mathbb{R} \times \mathbb{R}$ 的下述子集中的哪一个是 $\mathbb{R}$ 的两个子集的笛卡儿积。
\begin{enumerate}[(a)]
\item $\{(x, y) \mid x \text{ 为整数}\}$.
\item $\{(x, y) \mid 0 < y \leqslant 1\}$.
\item $\{(x, y) \mid y > x\}$.
\item $\{(x, y) \mid x \text{ 不是整数并且 } y \text{ 是一个整数}\}$.
\item $\{(x, y) \mid x^{2} + y^{2} < 1\}$.
\end{enumerate}
\end{exercise}

\end{document}